%%%%%%%%%%%%
%
% $Autor: Wings $
% $Datum: 2019-03-05 08:03:15Z $
% $Pfad: TemplateSensor $
% $Version: 4250 $
% !TeX spellcheck = en_GB/de_DE
% !TeX encoding = utf8
% !TeX root = filename 
% !TeX TXS-program:bibliography = txs:///biber
%
%%%%%%%%%%%%

% Structure
\chapter{Visual Studio Code}
	Visual Studio Code ist ein plattformübergreifender, quelloffener Code-Editor, der von Microsoft entwickelt wird und sich besonders für die Software- und Embedded-Entwicklung eignet. Durch seine modulare Architektur lässt sich die Entwicklungsumgebung mithilfe von Erweiterungen flexibel an unterschiedliche Programmiersprachen, Frameworks und Workflows anpassen.
	
\section{Installation Visual Studio Code für macOS}
	Die folgende \ac{vscode}-Installationsanleitung bezieht sich auf macOS.
	
	
	Die Installationsdatei für \ac{vscode} kann von der offiziellen Website heruntergeladen werden: \url{https://code.visualstudio.com/download}. Dabei ist darauf zu achten, die Version für macOS auszuwählen.
	
	
	\begin{center}		
		\includegraphics[width=0.7\linewidth]{VisualStudioCode/VSCWebsite}
		\captionof{figure}{Visual Studio Code Website}
		\label{fig:VisualStudioCodeWebsite}
	\end{center}
	
	Nachdem die Datei heruntergeladen wurde, öffnen wir den ,,Downloads'' - Ordner und kopieren ,,\ac{vscode}'' nach ,,Applications''. Dadurch erhalten wir Zugriff auf diese Applikation über unser Menü.
	
	\begin{center}		
		\includegraphics[width=0.7\linewidth]{VisualStudioCode/VSCDownload}
		\captionof{figure}{Heruntergeladene Applikation in Downloads Ordner}
		\label{fig:DownloadVSC}
	\end{center}
	
	\begin{center}		
		\includegraphics[width=0.7\linewidth]{VisualStudioCode/VSCCopy}
		\captionof{figure}{Kopieren der Applikation}
		\label{fig:CopyVSC}
	\end{center}
	
	An diesem Punkt lässt sich \ac{vscode} im Applikationsmenü finden und durch Anklicken starten.
	
	\begin{center}		
		\includegraphics[width=0.7\linewidth]{VisualStudioCode/VSCMenu}
		\captionof{figure}{Applikation Starten}
		\label{fig:VSCinMenu}
	\end{center}
	
	\begin{center}		
		\includegraphics[width=0.7\linewidth]{VisualStudioCode/VSCView}
		%\label{fig:ArduinoIDE Create Agent Installation}
		\captionof{figure}{Visual Studio Code Ansicht}
		\label{fig:VSCView}
	\end{center}
	
	Die Oberfläche des Programms kann je nach Version und Gerät leicht variieren.  